% !TEX program = lualatex
% Auto-generated from syllabus — 07: PCの設定2：OfficeとTeamsのインストール

\documentclass[aspectratio=43,professionalfonts,handout,10pt]{beamer}
\usepackage{beamer_template}

\newcommand{\LectureNum}{07}
\newcommand{\LectureTitle}{PCの設定2：OfficeとTeamsのインストール}
\newcommand{\CourseName}{情報リテラシー}
\newcommand{\TextPages}{-}

\begin{document}
% --- Slide 1: title ---

\begin{frame}{\CourseName\quad 第\LectureNum 回「\LectureTitle 」}
  {\small \TermName\quad \TeacherName\quad ／\quad 教科書 \TextPages}

  \vfill

  \begin{block}{今日の流れ}
    \begin{enumerate}
      \item 事前確認とインストール準備
      \item Microsoft 365（Office）のインストール
      \item Wordの動作確認
      \item Teamsのサインイン
    \end{enumerate}
  \end{block}
  \vfill

  \begin{block}{今日の目標}
    \begin{itemize}
      \item 各自のラップトップPCにOfficeとTeamsをインストールし，正常に動作することを確認する。
    \end{itemize}
  \end{block}
  \vfill

  \begin{exampleblock}{キーワード}
    \small \textbf{Microsoft 365} ／ \textbf{ライセンス} ／ \textbf{OneDrive} ／ \textbf{Teams}
  \end{exampleblock}
\end{frame}

% --- Slide 2: 作業の流れ ---

\begin{frame}{本日の作業の流れ}
  \begin{block}{ステップ}
    \begin{itemize}
      \item ① 事前確認とインストール準備（Wi-Fiオフ・有線LAN接続）
      \item ② Microsoft 365のインストール（portal.office.com）
      \item ③ Wordの動作確認（ファイル作成・保存）
      \item ④ Teamsのサインイン
    \end{itemize}
  \end{block}

  \vfill

  \begin{noticebox}[注意]
    インストール中は他の作業をしない。完了まで待つこと。
  \end{noticebox}
\end{frame}

% --- Slide 3: STEP 1 ---

\begin{frame}{STEP 1：事前確認とインストール準備}
  {\small Officeがインストール済みでないか確認し，ネットワークを準備する。}

  \vfill

  \begin{block}{手順}
    \begin{enumerate}
      \item スタートメニューで「Word」を検索し，インストール済みでないことを確認する
      \item Wi-Fiをオフにする
      \item USB有線LANアダプタを接続し，有線LANでインターネットに接続する
      \item ブラウザからインターネットにアクセスできることを確認する
    \end{enumerate}
  \end{block}

  \vfill

  \begin{noticebox}[注意]
    既にOfficeがインストールされている場合は担当教員・技術職員に相談すること。
  \end{noticebox}
\end{frame}

% --- Slide 4: STEP 2 ---

\begin{frame}{STEP 2：Microsoft 365のインストール}
  {\small 学校のアカウントでMicrosoft 365をインストールする。}

  \vfill

  \begin{block}{手順}
    \begin{enumerate}
      \item ブラウザで \kw{portal.office.com} を開く
      \item 学校のMicrosoftアカウントでサインインする
      \item 「Officeのインストール」→「Microsoft 365 Apps」をクリックする
      \item ダウンロードされたファイルを実行し，インストールが完了するまで待つ
    \end{enumerate}
  \end{block}

  \vfill

  \begin{noticebox}[注意]
    インストールには時間がかかる。途中で電源を切ったりブラウザを閉じたりしないこと。
  \end{noticebox}
\end{frame}

% --- Slide 5: STEP 3 ---

\begin{frame}{STEP 3：Wordの動作確認}
  {\small Wordを起動してファイルを作成・保存し，正常に動作することを確認する。}

  \vfill

  \begin{block}{手順}
    \begin{enumerate}
      \item スタートメニューからWordを起動する
      \item 「白紙の文書」を選択し，自分の名前を入力する
      \item 「名前を付けて保存」で \kw{OneDrive} に \kw{office\_ok\_学籍番号.docx} として保存する
      \item エクスプローラーでOneDriveフォルダを開き，保存したファイルを確認する
    \end{enumerate}
  \end{block}
\end{frame}

% --- Slide 6: STEP 4 ---

\begin{frame}{STEP 4：Teamsのサインイン}
  {\small Teamsに学校アカウントでサインインする。}

  \vfill

  \begin{block}{手順}
    \begin{enumerate}
      \item Teamsを起動し，学校のMicrosoftアカウントでサインインする
      \item サインイン後，チーム一覧やチャットが表示されることを確認する
    \end{enumerate}
  \end{block}
\end{frame}

% --- Slide 7: トラブル対応 ---

\begin{frame}{トラブル発生時の対処}
  \begin{noticebox}[よくあるトラブルと対処法]
    \begin{itemize}
      \item パスワードがわからない → 担当教員・技術職員に申し出る
      \item ダウンロードが始まらない → ブラウザを変えて試す（Edge・Chrome等）
      \item インストールが途中で止まる → 有線LAN接続を確認し，再試行する
      \item ライセンスエラーが出る → 一度サインアウトして再サインインする
    \end{itemize}
  \end{noticebox}

  \vfill

  \begin{exampleblock}{ポイント}
    解決できない場合は担当教員・技術職員に相談する。自己判断でアンインストール等をしないこと。
  \end{exampleblock}
\end{frame}

% --- チェックリスト ---

\begin{frame}{まとめ・チェックリスト}
  \begin{block}{自己チェック（できたら $\checkmark$ をつけよう）}
    \begin{itemize}
      \item[$\square$] Microsoft 365（Office）をインストールできた
      \item[$\square$] Wordを起動してファイルを作成・保存できた
      \item[$\square$] Teamsにサインインできた
    \end{itemize}
  \end{block}

\end{frame}

\end{document}
